%!TEX root = ../sztuthesis_main.tex
% 论文正文是主体，主体部分应从另页右页开始，每一章应另起页。一般由序号标题、文字叙述、图、表格和公式等五个部分构成。
\section{引言}
\subsection{研究背景与意义}

近年来，AI 模型技术发展迅猛，尤其是大语言模型（Large Language Model, LLM）在自然语言处理领域的突破性进展和落地应用。
人们逐渐抛弃传统的搜索引擎转而直接使用 AI 软件获取信息的场景越来越多。
但受限于模型训练时的输入数据，模型对训练数据之外的信息缺乏感知能力，这也就导致了 LLM 模型、在某些特定领域的应用受限，无法满足用户的个性化需求。

\subsection{课题基础及国内外研究进展}
为了解决模型无法得知外部数据的问题，近年来 Retrieval Augmented Generation（RAG）技术被提出并广泛应用。
RAG 技术根据用户的 prompt 来检索相关的外部知识库文档，并在推理中一起提交给模型作为输入，让模型能够得知新的信息。
RAG 技术确实有效，但仍然有一些问题，比如：难以根据文档被创建的时间关系逻辑来推理，每次推理都需要重新检索，检索结果可解释性低等。
RAG 技术缺乏对知识的统一组织和长期维护能力。

\subsection{主要研究工作}
本文围绕着高效利用个人知识库来拓宽 LLM 模型的应用边界问题展开研究，主要工作包括：

(1) 处理用户过往的笔记/博客/对话内容，提取出事实（Fact），基于此描述用户的背景信息。

(2) 根据 Fact 推理出用户的兴趣和需求，基于此推荐给用户一些能力提升的资料，并能规划和持续追踪用户的学习进度。

(3) 记录用户的知识库变化日志，基于此分析用户最近的活动状态。

\subsection{论文组织结构}

全文内容共六章，具体内容组织如下：

第一章为绪论。

第二章为图片插入示例。

第三章为表格插入示例。

第四章为公式插入示例。

第五章为参考文献插入示例。

第六章总结与展望，总结了本文的主要工作，展望了下一阶段的研究方向。

% \newpage    % 两个章节之间分页，不想分的话可注释掉

\section{图像布局}
\label{sec.figure}

\emph{学校对图片只有小标题要求，没有进一步的子图要求，我们按科技论文常规排版来}

\subsection{单图布局}

\lipsum

\emph{单图布局如图~\ref{F.sztu_single} 所示。}

\begin{figure}[hbt]
  \centering
  \includegraphics[width=0.5\textwidth]{sztu.png}
  \caption{单图布局示例}
  \label{F.sztu_single}
\end{figure}

\subsection{位置固定（\texttt{[H]}，非浮动）}

学校《撰写规范》对\textbf{图题字号、按章编号、图题在图下}等作了要求，\textbf{并未规定}插图必须采用浮动体；因此使用 \texttt{figure[H]} 将插图\textbf{严格固定在源码位置}（不漂到上下页）\textbf{不违反规范}。但 \texttt{[H]} 会阻止 \LaTeX{} 按版面优化重排插图，大图或连续多图时可能出现\textbf{分页难看、垂直空白过大}等问题，\textbf{宜按需少量使用}；一般情况仍推荐 \texttt{htbp}、\texttt{!htb} 等浮动策略。

以下示例中，插图应出现在\textbf{两段文字之间}，而不会「顶开」后一段到上一页或漂到后一页的浮动区。

第一段占位文字：\lipsum[1]

\begin{figure}[H]
  \centering
  \includegraphics[width=0.35\textwidth]{sztu.png}
  \caption{非浮动单图示例（\texttt{[H]}）}
  \label{F.here_fixed}
\end{figure}

第二段占位文字：\lipsum[2]

图~\ref{F.here_fixed} 使用 \pkg{float} 宏包提供的 \texttt{H} 说明符；本模板已在 \cls{SZTUthesis.cls} 中加载 \pkg{float}，正文中直接写 \verb|\begin{figure}[H]| 即可。

\subsection{横排布局}

\emph{横排布局如图~\ref{F.sztu_row} 所示。}

\begin{figure}[!htb]
  \centering
  \begin{subfigure}[t]{0.24\linewidth}
    \begin{minipage}[b]{1\linewidth}
      \includegraphics[width=1\linewidth]{sztu.png}
      \caption{可以增加描述}
    \end{minipage}
  \end{subfigure}
  \begin{subfigure}[t]{0.24\linewidth}
    \begin{minipage}[b]{1\linewidth}
      \includegraphics[width=1\linewidth]{sztu.png}
      \caption{}
    \end{minipage}
  \end{subfigure}
  \begin{subfigure}[t]{0.24\linewidth}
    \begin{minipage}[b]{1\linewidth}
      \includegraphics[width=1\linewidth]{sztu.png}
      \caption{}
    \end{minipage}
  \end{subfigure}
  \begin{subfigure}[t]{0.24\linewidth}
    \begin{minipage}[b]{1\linewidth}
      \includegraphics[width=1\linewidth]{sztu.png}
      \caption{}
    \end{minipage}
  \end{subfigure}
  \caption{横排布局示例}
  \label{F.sztu_row}
\end{figure}

\lipsum

\subsection{竖排布局}
\emph{竖排布局如图\ref{F.sztu_col}所示。}

\begin{figure}[!htb]
  \centering
  \begin{subfigure}[t]{0.15\linewidth}
    \begin{minipage}[b]{1\linewidth}
      \includegraphics[width=1\linewidth]{sztu.png}
      \caption{}
    \end{minipage}
  \end{subfigure}\\
  \begin{subfigure}[t]{0.15\linewidth}
    \begin{minipage}[b]{1\linewidth}
      \includegraphics[width=1\linewidth]{sztu.png}
      \caption{}
    \end{minipage}
  \end{subfigure}
  \caption{竖排布局示例}
  \label{F.sztu_col}
\end{figure}

\lipsum

\subsection{竖排多图横排布局}

\begin{figure}[!htb]
  \centering
  \begin{subfigure}[t]{0.13\linewidth}
    \begin{minipage}[b]{1\linewidth}
      \includegraphics[width=1\linewidth]{sztu.png} \vspace{-1ex} \vfill
      \includegraphics[width=1\linewidth]{sztu.png}
    \end{minipage}
    \caption{}
  \end{subfigure}
  \begin{subfigure}[t]{0.13\linewidth}
    \begin{minipage}[b]{1\linewidth}
      \includegraphics[width=1\linewidth]{sztu.png} \vspace{-1ex} \vfill
      \includegraphics[width=1\linewidth]{sztu.png}
    \end{minipage}
    \caption{}
  \end{subfigure}
  \caption{竖排多图横排布局}
  \label{F.sztu_col_row}
\end{figure}

\emph{竖排多图横排布局如图~\ref{F.sztu_col_row} 所示。注意看(a)、(b)编号与图关系。}

\subsection{横排多图竖排布局}

\lipsum

\begin{figure}[!htb]
  \centering
  \begin{subfigure}[t]{0.3\linewidth}
    \begin{minipage}[b]{1\linewidth}
      \includegraphics[width=0.45\linewidth]{sztu.png}
      \includegraphics[width=0.45\linewidth]{sztu.png}
    \end{minipage}
    \caption{}
  \end{subfigure}\\
  \begin{subfigure}[t]{0.3\linewidth}
    \begin{minipage}[b]{1\linewidth}
      \includegraphics[width=0.45\linewidth]{sztu.png}
      \includegraphics[width=0.45\linewidth]{sztu.png}
    \end{minipage}
    \caption{}
  \end{subfigure}
  \caption{横排多图竖排布局}
  \label{F.sztu_row_col}
\end{figure}

\emph{横排多图竖排布局如图~\ref{F.sztu_row_col} 所示。注意看(a)、(b)编号与图关系。}

\subsection{本章小结}
本章示例图片布局。

这里再测试一下不同章节的公式编号
\begin{equation}
  p_{i} = \frac{e^{-\varepsilon_{i}/kT}}{\sum_{j=1}^{M}e^{-\varepsilon_{j}/kT}}
\end{equation}

\newpage    % 两个章节之间分页，不想分的话可注释掉

\section{表格插入示例}

\begin{table}[htb]
  \centering
  \caption{学校文件里对表格的要求不是很高，不过按照学术论文的一般规范，表格为三线表。}
  \label{T.example}
  \begin{tabular}{llllll}
    \hline
    & A  & B  & C  & D  & E \\
    \hline
    1   & 212 & 414 & 4     & 23 & fgw  \\
    2   & 212 & 414 & v     & 23 & fgw  \\
    3   & 212 & 414 & vfwe    & 23 & 长一些的内容  \\
    4   & 212 & 414 & 4fwe    & 23 & 嗯  \\
    5   & af2 & 4vx & 4     & 23 & 长一些的内容  \\
    6   & af2 & 4vx & 4     & 23 & fgw  \\
    7   & 212 & 414 & 4     & 23 & fgw  \\

    \hline{}
  \end{tabular}
\end{table}

\emph{表格如表~\ref{T.example} 所示，\LaTeX 表格技巧很多，这里不再详细介绍。}

\lipsum

\newpage    % 两个章节之间分页，不想分的话可注释掉

\section{公式插入示例}

\lipsum

\emph{公式插入示例如公式~\eqref{E.example} 所示。}
\begin{equation}
  \gamma_x=
  \begin{cases}
    0, & \text{if $|x| \leq \delta$} \\
    x, & \text{otherwise}
  \end{cases}
  \label{E.example}
\end{equation}

\newpage    % 两个章节之间分页，不想分的话可注释掉

\section{参考文献插入示例}

\LaTeX \cite{lamport1994latex}插入参考文献最方便的方式是使用 \env{bibliography}\cite{pritchard1969statistical}。

大多数出版商的论文页面都会有导出 \format{bib} 格式参考文献的链接，把每个文献的 \format{bib} 放入 \bib{thesis-references.bib}，然后用 \oper{bibkey} 即可插入参考文献。

\lipsum

\newpage    % 两个章节之间分页，不想分的话可注释掉

\section{总结与展望}

\noindent{纯数字编号}
\begin{enumerate}
  \item XXXXXXXXXX
    \label{item1}
  \item XXXXXXXXXX
  \item XXXXXXXXXX
\end{enumerate}
罗马编号
\begin{enumerate}[label=(\roman*)]
  \item XXXXXXXXXX
    \label{item2}
  \item XXXXXXXXXX
  \item XXXXXXXXXX
\end{enumerate}
括号编号
\begin{enumerate}[label=(\arabic*)]
  \item XXXXXXXXXX
    \label{item3}
  \item XXXXXXXXXX
  \item XXXXXXXXXX
\end{enumerate}
半括号编号
\begin{enumerate}[label=\arabic*)]
\item XXXXXXXXXX
  \label{item4}
\item XXXXXXXXXX
\item XXXXXXXXXX
\end{enumerate}
小字母编号
\begin{enumerate}[label=\alph*)]
\item XXXXXXXXXX
\label{item5}
\item XXXXXXXXXX
\item XXXXXXXXXX
\end{enumerate}

引用测试,正如~\ref{item1}、\ref{item2}、\ref{item3}、\ref{item4}、\ref{item5} 所示

\subsection{工作展望}
手动编号 %（不推荐,无法被交叉引用）

本课题针对XX，鉴于XXX，对XX进行了提高，但是XXX，所以有如下XX：

（1）目前XX虽然XX，但是XX仍然XX，所以XX仍然是一个值得XX的问题。

（2）随着XX，XX具有XX的问题，仍值得进一步XX。

（3）本课题在XX有了XX，但是XX的XX还存在XX，所以XX。

\newpage
